\cleardoublepage
\chapter*{Abstract}
\markboth{Abstract}{}
\addcontentsline{toc}{chapter}{Abstract}

Breast cancer remains one of the deadliest diseases affecting women
worldwide, making accurate diagnosis and treatment planning crucial
for improving patient outcomes. This thesis presents a novel approach
to Image Semantic Segmentation (ISS) in medical imaging, with a
particular focus on histopathology images of breast cancer tissue.
The work addresses the critical challenge of handling multiple
annotators' labels in crowdsourcing scenarios, where annotator
performance varies across the input space and image quality. The
proposed solution introduces a novel loss function that implicitly
infers optimal segmentation and assesses labeler reliability without
requiring explicit performance inputs. This is achieved through an
intermediate reliability map that enables the model to prioritize
information from reliable labelers while disregarding noisy labels.

The methodology diverges from conventional CNN-based segmentation
models by eliminating the need for direct supervision of labeler
performance, thereby enhancing its generalizability and adaptability
across diverse datasets and labeler cohorts. The approach combines
the strengths of U-shaped deep learning architectures with a
sophisticated loss function that models annotator reliability across
the input space. The model's performance is evaluated on both
synthetic datasets and real-world histopathology images,
demonstrating superior results compared to state-of-the-art
approaches in terms of segmentation accuracy and uncertainty quantification.

The research makes significant contributions to the field by
providing a robust solution for handling inconsistent annotations in
medical image segmentation, particularly valuable in histopathology
where expert annotations are expensive and time-consuming to obtain.
The proposed approach shows promising results in reducing annotation
costs while maintaining high segmentation accuracy, making it
particularly relevant for clinical applications where precise
segmentation of tissue structures is crucial for diagnosis and
treatment planning.

\textbf{\small Keywords:} Crowdsourcing, Multiple Annotators,
Histopathology, Breast Cancer, Semantic Image Segmentation, Gaussian
Processes, Deep Learning, Medical Image Analysis
